% 06_body.tex — Day2 산출물 1/5 (⑥) 범위 기술서와 WBS, WBS 사전 (빈 템플릿 / 완성 예시 공용 본문)
% 드라이버: 06_범위기술서WBS_빈템플릿.tex (\wbblanktrue) / 06_범위기술서WBS_완성예시.tex (\wbblankfalse)
% 내용 출처: PM트랙/Day2_워크북.md §1.3(항목 가이드) §1.4(빈 템플릿) §1.5(온담 P1 완성 예시본)
\documentclass[10pt,a4paper]{article}
\usepackage[a4paper,margin=16mm,top=14mm,bottom=20mm]{geometry}
\def\DocNum{1}
\def\DocTitle{범위 기술서와 WBS, WBS 사전}
\def\DocSub{⑥ Scope Statement, WBS \& WBS Dictionary}
\def\DocVariant{\B{빈 템플릿}{온담 P1 완성 예시}}
\usepackage{wbstyle}
\begin{document}

\docheader
 {\Bg{[정식 명칭과 코드]}{차세대 주문·재고 통합 플랫폼 구축 (ONE ONDAM P1)}}
 {\Bg{[v0.x / 기준선 상태]}{v0.9 / 초안 — 2026-03-27 (S6 종료), G1 승인 예정}}
 {\Bg{[PM 실명]}{P1 PM (발주사 IT본부)\rd{7mm}}}
 {\Bg{[스폰서 실명]}{정미란 CFO (G1 2026-04-30)\rd{7mm}}}

\begin{center}\bfseries\large 범위 기술서 (Scope Statement)\end{center}
% ------------------------------------------------------------------ 1
\secbar{1. 문서 정보와 기준선 상태}
\guide{프로젝트명·코드, 문서 버전, 상위 문서(헌장 vN), 기준선 상태(초안 / G1 승인 / 변경 n차), 작성자·검토자·승인자, 관련 문서(RTM 버전, FP 기준선 확인서 번호). 기준선 상태 칸에는 다음 승인 시점을 함께 적는다. 범위 기술서 버전과 FP 기준선 버전이 따로 놀면 변경 시 어느 문서를 고쳤는지 아무도 모른다.}
{\footnotesize
\begin{tabularx}{\linewidth}{|L{30mm}|Y|L{22mm}|Y|}\hline
\hd{항목} & \hd{내용} & \hd{항목} & \hd{내용} \\ \hline
\B{\rd{9mm}}{}프로젝트명 / 코드 & \Bg{[정식 명칭 / 코드]}{차세대 주문·재고 통합 플랫폼 구축 / ONE ONDAM P1} & 상위 문서 & \Bg{[헌장 vN, 일자]}{프로젝트 헌장 v1.0 (2026-01-05)} \\ \hline
\B{\rd{9mm}}{}버전 / 기준선 상태 & \Bg{[v0.x / 초안·G1 승인·변경 n차]}{v0.9 / 초안 — 설계 스프린트 S6 종료(2026-03-27) 시점, G1 승인 예정} & 관련 문서 & \Bg{[RTM vN, FP 기준선 확인서 번호]}{RTM v0.9, FP 산정서 v0.9, FP 기준선 확인서 BL-01(G1 서명 예정)} \\ \hline
\B{\rd{9mm}}{}작성자 / 검토자 / 승인자 & \Bg{[실명 / 실명 / 실명]}{P1 PM / 이재현 수행사 PM, 박세연 IT본부장, 오정근 영업본부장 / 정미란 CFO} & 승인 예정 & \Bg{[게이트, 날짜]}{G1 2026-04-30 (프로그램 이사회)} \\ \hline
\end{tabularx}}

% ------------------------------------------------------------------ 2
\secbar[70mm]{2. 포함 범위 — 산출물 단위}
\guide{헌장 §5의 포함 산출물을 산출물 단위로 다시 쓰되, 각 산출물에 “무엇이 포함되는가”를 3~6줄로 구체화한다. 각 줄이 WBS 레벨 3 작업패키지 하나 이상에 대응해야 한다. 헌장 §13에서 “이연 가능”으로 표시된 것은 여기서도 “이연 가능(릴리스 Rn)”으로 표시. 헌장 문장 복사(범위 기술서가 존재할 이유가 없다)도, SRS 570건 나열(RTM의 일)도 아니다 — 30~50줄이 적정.}
\begin{fieldbox}
\textbf{산출물 ①} \Bg{[명칭]}{신규 클라우드 POS·매장 운영 시스템}
\ifwbblank\lines{5}{8mm}\else
\begin{itemize}
\item 판매·결제·반품·할인·멤버십 적립의 매장 거래 처리, 중앙 재고 실시간 연결(로컬 DB 폐지)
\item 매장 운영: 개점·마감, 재고 실사 입력, 매장 발주 연계, 폐기 등록
\item 결제 연동: 은행·PG·밴 인터페이스 7본 재연결(문서 부재 2본 복원 포함)
\item 직영 227점 단말 477대 신규 조달·설치, 가맹 388점 기존 단말에 SW 배포 패키지 제공
\item 파일럿 12개점(P3 지정) 컷오버 전 배포, 615개 매장 펌웨어 일괄배포 패키지와 롤백 절차
\end{itemize}\fi
\end{fieldbox}
\begin{fieldbox}
\textbf{산출물 ②} \Bg{[명칭]}{주문 허브(Order Hub)}
\ifwbblank\lines{5}{8mm}\else
\begin{itemize}
\item 자사앱·배달3사·매장 POS·가맹 발주·B2B EDI 5개 채널의 주문을 하나의 주문 원장에 기록
\item 배달3사 인터페이스 3본 재연결, 자사앱 API 규격 제공(앱측 구현은 앱 벤더)
\item B2B 거래처 EDI(편의점 3사·대형식품사 5사·원료유통 260개사) 주문 접수
\item 옴니채널 기본 시나리오 2종: 앱 주문 매장 픽업, 매장 간 재고 이동 요청
\item \textbf{이연 가능(릴리스 R5)}: 매장 간 재고 이동 자동 제안, 픽업 예약 시간 지정, 주문 BI 리포트 2종 (헌장 §13)
\end{itemize}\fi
\end{fieldbox}
\begin{fieldbox}
\textbf{산출물 ③} \Bg{[명칭]}{통합 재고 관리}
\ifwbblank\lines{5}{8mm}\else
\begin{itemize}
\item 매장·이천 물류센터·안성·이천 공장 출하 재고의 단일 재고 원장, 반영 지연 15분 이내
\item 3PL 대성로지스 WMS 준실시간 연동 명세서와 P1측 API, 폴백 야간 배치 2회(21:00·03:00)
\item P2 설비 인터페이스(15억, 별도 계약)의 P1측 데이터 규격과 수신 구현 — 설비측 구현·시험은 P2
\item ㈜온담푸드시스템 공장 출하 재고 연동(별도 법인 계약 구조 확정 후)
\item 재고 정확도 리포트 — 순환실사 SKU 69\%와 전체 SKU 100\% 두 분모 병기(R-12)
\end{itemize}\fi
\end{fieldbox}
\begin{fieldbox}
\textbf{산출물 ④} \Bg{[명칭]}{가맹 발주 포털}
\ifwbblank\lines{5}{8mm}\else
\begin{itemize}
\item 발주·마감·출고 확인·배송 추적 단일 화면, 발주 리드타임 12시간 설계
\item 전화·팩스·엑셀 메일 39\% 흡수를 위한 엑셀 업로드·모바일 발주
\item 388점 계정 발급·온보딩, 협의회 참여 프로토타입 시연(2026-05), 선행 릴리스(2026-Q4, O-02)
\item 가맹 발주 웹 ↔ SAP 인터페이스 4본 대체(문서 부재 2본 복원 포함)
\end{itemize}\fi
\end{fieldbox}
\begin{fieldbox}
\textbf{산출물 ⑤} \Bg{[명칭]}{통합 계층·데이터 전환·공통 기반}
\ifwbblank\lines{5}{8mm}\else
\begin{itemize}
\item API 게이트웨이, 상품·고객 MDM, SAP 애드온(ECC 6.0 기준, 2027-12 이후 마이그레이션 고려 분리 구조)
\item 인터페이스 47본 재구축·재연결, 설계문서 47본 전량(문서 부재 13본 복원)
\item 데이터 전환: 1,842 테이블·7.4TB 중 5년 이내 이력 이관, 5년 초과분은 아카이브(신규 가정, §5), 마스터 6종 정본(G2), 리허설 3회, 본 전환
\item 공통 기반: 클라우드(선약정), DR·백업, 망분리·보안장비, APM·모니터링
\item 개인정보 처리항목 142개·위탁사 19개·재위탁 4건 설계 반영, 전자금융 요건 반영
\end{itemize}\fi
\end{fieldbox}

% ------------------------------------------------------------------ 3 (가로)
\landscapeon
\secbar{3. 제외 범위와 경계 조건}
\guide{헌장 §5 제외 항목을 담당과 함께 옮기고, 각 항목에 “P1이 그 항목과 맞닿는 지점에서 무엇까지는 한다”를 적는다. 경계 조건은 P1이 인도하는 것 / 상대가 인도하는 것 / 상대가 인도하지 않을 때의 폴백 세 칸. \textbf{폴백 칸이 비어 있는 경계 조건은 리스크 등록부에 위협으로 등록되어 있어야 한다.} 제외만 적고 경계에서 P1이 하는 일을 적지 않으면 “3PL 연동을 안 한다”와 “3PL 연동 명세를 쓴다”가 한 문서 안에서 충돌한다.}
{\footnotesize
\begin{longtable}{|L{44mm}|L{34mm}|L{58mm}|L{40mm}|L{62mm}|}\hline
\hd{제외 항목} & \hd{담당 (프로젝트·부서·실명)} & \hd{P1이 여기까지는 함} & \hd{P1이 이것은 안 함} & \hd{폴백 (상대가 인도하지 않을 때) 또는 R-nn} \\ \hline \endhead
\ifwbblank
\rd{13mm}\g{[제외 항목 1]} & \g{[프로젝트·부서·실명]} & \g{[P1이 인도하는 것]} & \g{[상대가 인도하는 것]} & \g{[폴백 또는 R-nn]} \\ \hline
\rd{13mm} & & & & \\ \hline
\rd{13mm} & & & & \\ \hline
\rd{13mm} & & & & \\ \hline
\rd{13mm} & & & & \\ \hline
\rd{13mm} & & & & \\ \hline
\rd{13mm} & & & & \\ \hline
\rd{13mm} & & & & \\ \hline
\else
SAP ECC 6.0 교체·업그레이드 & IT본부 별도 과제(박세연) & 애드온·API를 ECC 기준으로 구축, 마이그레이션 고려 분리 구조 & ECC 자체 변경, S/4 전환 & — (2027-12까지 ECC 유지) \\ \hline
이천 물류센터 설비·자동화·WMS & P2 (한동석, 타케다 PM) & 재고 원장 기준 데이터 규격 T0(2026-04-24) 전 전달, P1측 수신 구현·통합시험 & 설비측 구현, 설비 시험 일정, WMS 자체 & 컷오버 물류 부분 2차 분리 — 프로그램 이사회 결정 (R-01) \\ \hline
3PL 대성로지스 계약 변경 협상 & 물류본부 (한동석) & 준실시간 연동 명세서, 부하 산정서, P1측 API 구현 & 협상, 계약 변경 서명 & 야간 배치 2회 폴백, 재고 정확도 목표는 매장 재고 한정 (R-03) \\ \hline
매장 직원·가맹점주 교육, 슈퍼유저, 가맹 계약 부속서 개정 & P3 (오정근) & UAT 참여자 선정 지원, 교육용 환경 제공, 데이터 용도 명세서 & 교육 실시, 388점 동의 협상 & 가맹 SW 배포 범위를 동의 매장으로 축소 (R-02) \\ \hline
개인정보 보호법 대응 완료 판정, 전자금융업 등록 신청, ISMS-P 취득 & P4 (최영주) & P4가 확정한 아키텍처 요건 설계 반영, 142개 처리항목 설계 검증 자료 & 요건 해석, 규제기관 대응 & 규제 버퍼 22영업일 소진 시 P4 요건 흡수 — 흡수 결정은 프로그램 계층 \\ \hline
운영 조직 편성, SLA 체결, 서비스데스크 & P5 (박준영) & 운영 이관 기술 문서(12종 중 시스템 기술 문서), SAC 초안 공동 작성 & 운영 절차 문서, 조직 편성 & 운영 절차 미완 시 기술 문서에 임시 운영 절차 부록 (R-11) \\ \hline
가맹점 POS 단말 하드웨어 교체 & 가맹 계약 협상 결과에 따름 (오정근) & 388점 단말 기종 전수 조사(2026-04), 비호환 기종 목록·대응안 & 교체 비용 부담 & 비호환 매장은 컷오버 후 수기 병행 + 2차 배포 (R-13) \\ \hline
자사앱 UI/UX 개편 & 영업본부·앱 벤더 (오정근) & 앱 주문 연동 API 규격 제공, 연동 시험 환경 & 앱 화면·기능 개편, 앱측 구현 & 앱측 미구현 시 앱 주문은 기존 경로 유지, B-05 지연 (R-08) \\ \hline
\fi
\end{longtable}}

% ------------------------------------------------------------------ 4 (가로)
\secbar{4. 산출물 인수 기준 (Acceptance Criteria)}
\guide{헌장 §6 주요 산출물 각각에 “무엇을 확인하면 인수인가”를 참·거짓으로 판정 가능한 문장 2~4개로. 인수자·인수 시점은 헌장에서 가져오고, 여기서 새로 붙이는 것은 \textbf{판정 기준과 판정 방법}이다. 숫자는 헌장 §3 성공 기준·§4 상위 요구를 인용(가용성 99.5\%, 15분, 12시간). “정상 동작 확인”, “사용자 만족”은 판정 불가. 소스·형상 인수 기준(클린 빌드 재현)을 빠뜨리지 말 것.}
{\footnotesize
\begin{longtable}{|L{36mm}|L{104mm}|L{46mm}|L{28mm}|L{34mm}|}\hline
\hd{산출물} & \hd{인수 기준 (참·거짓 판정 가능, 2~4개)} & \hd{판정 방법 (시험 ID·검수·감리·문서)} & \hd{인수자} & \hd{인수 시점} \\ \hline \endhead
\ifwbblank
\rd{14mm}\g{[산출물 1]} & \g{① [\ ] ② [\ ] ③ [\ ]} & \g{[시험 ID / 검수 조서 / 감리 / 문서 확인]} & \g{[실명]} & \g{[게이트·날짜]} \\ \hline
\rd{14mm} & & & & \\ \hline
\rd{14mm} & & & & \\ \hline
\rd{14mm} & & & & \\ \hline
\rd{14mm} & & & & \\ \hline
\rd{14mm} & & & & \\ \hline
\rd{14mm} & & & & \\ \hline
\rd{14mm}소스·형상 인수 & \g{① 클린 빌드 재현 [\ ]\% ② [3자 컴포넌트·라이선스] ③ [절차서]} & & & \\ \hline
\else
신규 클라우드 POS·매장 운영 시스템 & ① 직영 477대 설치·거래 처리 성공 100\% ② 가맹 배포 패키지가 호환 기종 전 유형에서 설치·거래 성공 ③ 성수기 창 밖 배포 잠금 동작 ④ 롤백 리허설 2회 기록 & 통합시험 TC-POS-001~180, 설치 확인서, 롤백 리허설 기록 RH-01~02 & 오정근, 박세연 & G3 2027-02-19, 컷오버 \\ \hline
주문 허브 & ① 5개 채널 주문이 채널 식별자와 함께 단일 원장에 기록 ② 픽업·매장 간 이동 기본 시나리오 2종 UAT 통과 ③ 배달3사 3본·B2B EDI 접수 통합시험 통과 & 통합시험 TC-ORD-001~085, UAT-ORD-01~05 & 오정근 & G3 (고도화 R5는 G4) \\ \hline
통합 재고 관리 & ① 3개 재고원 단일 원장 조회, 반영 지연 15분 이내 ② 3PL 준실시간 또는 폴백 자동 전환 시험 통과 ③ 재고 정확도 리포트 두 분모 병기 ④ P2 인터페이스 수신 시험 통과(설비측 MC 후) & 통합시험 TC-INV-001~040, 성능시험 & 한동석, 박세연 & G3 \\ \hline
가맹 발주 포털 & ① 388점 계정 발급 100\% ② 발주→출고 확정 12시간 설계 검증(시험 데이터) ③ 협의회 지정 점주 5명 사용성 검수 통과 & UAT TC-FRN-001~060, 사용성 검수 조서 & 오정근, 서정필(검수 참여) & 선행 릴리스 2026-Q4, G3 \\ \hline
통합 계층·데이터 전환 결과 & ① API GW·MDM·SAP 애드온 가동, 인터페이스 47본 통합시험 통과 ② 전환 검증 6종(건수·합계·키 정합·참조 무결성·표본 대사·업무 검증) 통과 ③ 리허설 3회차 34시간 이내 & 통합시험 TC-INT-001~094, 전환 검증 보고서 & 박세연 & 리허설 3회차 2027-01-29 → 컷오버 \\ \hline
인터페이스 설계문서 47본 & ① 47본 전량, 문서 부재 13본 포함 ② 표준 양식(연동 주체·주기·항목·오류 처리) 충족 & 감리 문서 검토 & 박세연 & G2 2026-11-27 \\ \hline
소스·형상 인수 & ① 발주사 저장소에서 클린 빌드 재현 100\% ② 3자 컴포넌트 목록(SBOM)·라이선스 제출 ③ 빌드·배포 절차서 & 형상관리 담당 재현 시험 & 박세연, 발주사 형상관리 담당 & G4 2027-05-29 \\ \hline
운영 이관 기술 문서 & ① SAC 기술 항목 전부 대응 문서 존재 ② 장애 대응 절차 3건 이상 실제 수행 기록 & SAC 점검표 & 박준영 & G4 \\ \hline
\fi
\end{longtable}}
\landscapeoff

% ------------------------------------------------------------------ 5
\secbar[70mm]{5. 범위 관련 가정과 제약}
\guide{헌장 §9의 가정·제약 중 범위에 직접 영향을 주는 것만 골라 옮기고, 작성 중 새로 발견한 가정은 “[신규]”로 표시해 헌장 v1.1 수정 요청으로 잇는다. 가정에는 확인 시점·확인 책임자, 제약에는 출처를 붙인다. 범위와 무관한 헌장 가정까지 전부 복사하지 말 것.}
\begin{fieldbox}
\textbf{가정} \guide{[내용] (확인: 책임자, 시점) — 출처 / [신규] 표시}
\ifwbblank\lines{6}{8.5mm}\else
\begin{itemize}
\item 3PL 계약 변경이 2026-06-30까지 타결된다 (확인: 한동석, 미타결 시 폴백) — 헌장 가정 1
\item 가맹 부속서 개정이 2026-08-14까지 388점 동의를 확보한다 (확인: 오정근) — 헌장 가정 2
\item 앱 벤더가 P1 API 규격에 맞춰 2026-12까지 연동을 완료한다 (확인: 오정근·박세연) — 헌장 가정 4
\item P2 MC 2027-02-12 달성, 물류 인터페이스 통합시험 G3 전 5영업일 확보 (확인: 프로그램 PMO) — 헌장 가정 5
\item \textbf{[신규]} 파일럿 12개점은 컷오버 전(2027-01) OD-POS와 신 POS를 병행하며, 그 기간의 임시 재고 동기화 연동(1.1.7)이 필요하다 (확인: P3 PM·박세연, 2026-06) — 헌장 v1.1 범위 수정 요청 대상
\item \textbf{[신규]} 데이터 전환 대상 이력 7.4TB 중 5년 초과분(약 2.6TB 추정 [추가 설정])은 아카이브 저장소에 보관하고 신 플랫폼에는 이관하지 않는다 (확인: 박세연·재경팀장 — 세무·감사 보존 요건, 2026-04) — 원가 기준선 조정 근거(⑨ 항목 3)
\end{itemize}\fi
\end{fieldbox}
\begin{fieldbox}
\textbf{제약} \guide{[내용] (출처) — 헌장 제약 n}
\ifwbblank\lines{4}{8.5mm}\else
\begin{itemize}
\item 컷오버 전까지 OD-POS 2.1·SAP ECC 6.0 병행 운영 (출처: IT본부) — 헌장 제약 3
\item ㈜온담푸드시스템 별도 법인 — 공장 연동은 별도 계약·개인정보 처리 주체 (출처: 법무) — 헌장 제약 7
\item 수행사 계약 FP 12,400점 고정가, 요구사항 동결 월 1회, FP 누적 변동 ±3\% 시 기준선 재확인 (출처: 계약서 부속서 1) — 헌장 제약 8
\item 대성로지스 WMS 커스터마이징 불가 (출처: 3PL 계약) — 헌장 제약 5
\end{itemize}\fi
\end{fieldbox}

% ------------------------------------------------------------------ 6
\secbar[70mm]{6. 범위 변경 규칙}
\guide{범위 기준선을 바꾸는 절차와 트리거. Day1 거버넌스(변경협의체 → CCB → 계약 변경)를 인용하되 범위 특유의 트리거 — 월 요구사항 동결, FP 누적 변동 ±3\%, 릴리스 이연 — 를 명시. 허용범위의 단위(건 / FP / \%)는 허용범위 격자(⑩ 항목 7)의 범위 축과 같아야 한다. “변경관리 계획 참조”로 끝내지 말 것.}
\begin{fieldbox}
\ifwbblank
\g{허용범위 내 (±\hspace{10mm}건, FP ±\hspace{10mm}점): 승인자 \hrulefill}\par\vspace{5mm}
\g{초과 시: \hrulefill}\par\vspace{5mm}
\g{동결 주기·접수 마감: \hrulefill}\par\vspace{5mm}
\g{FP 누적 ±\hspace{8mm}\%: 기준선 재확인 — 서명자 \hrulefill}\par\vspace{5mm}
\g{릴리스 이연: 승인 회의체 \hspace{30mm} 인용 문서 \hrulefill}\par\vspace{5mm}
\g{Must 상한: \hrulefill}\par
\else
\begin{itemize}
\item 요구사항은 매월 마지막 수요일에 동결한다. 동결 2주 전 수요일에 접수를 마감하며(R-07 대응), 마감 후 접수분은 다음 달 동결로 넘긴다.
\item MoSCoW 분류에서 Must는 릴리스별 FP의 60\% 이하로 유지한다. 초과 시 Senior User(오정근)가 Should로 내릴 항목을 지정한다.
\item 허용범위 내 변경(월 순변동 ±35건 이하, FP 누적 ±372점 이하)은 P1 PM이 승인하고 CCB에 사후 보고한다. 초과 시 CCB 승인, 변경협의체 합의 선행.
\item FP 누적 변동이 ±3\%(±372점)를 넘는 순간 기준선 재확인(P1 PM·이재현·형상관리 담당 3인 서명)을 실시한다. 분기 1회 정기 재확인과 별개다.
\item 릴리스 이연(R5로 이연 가능 항목의 실제 이연 결정)은 헌장 §13을 인용해 CCB가 승인한다. 이연 항목은 범위 기준선에서 빠지는 것이 아니라 인도 시점이 바뀌는 것이며, FP 기준선은 유지된다.
\item 범위 기술서·WBS·RTM·FP 기준선 확인서는 하나의 범위 기준선이다. 어느 하나를 바꾸면 넷 모두의 버전이 올라간다.
\end{itemize}\fi
\end{fieldbox}

\clearpage
\begin{center}\bfseries\large WBS (Work Breakdown Structure)\end{center}
% ------------------------------------------------------------------ 7
\secbar[100mm]{7. WBS 구조 규칙 (레벨·코드·분해 기준)}
\guide{레벨의 의미(L1 프로젝트 / L2 산출물 / L3 작업패키지), 코드 체계(1.1.1), 분해 기준(산출물 지향, 100\% 규칙, 작업패키지 크기 상하한), 프로젝트 관리 작업의 위치를 5~7줄로. \textbf{L2에 “프로젝트 관리” 노드를 반드시 둔다} — PMO·감리·통합시험·컷오버 관리가 갈 곳. 활동(“설계하기·개발하기·시험하기”)으로 분해하면 산출물 하나의 원가가 세 노드에 흩어지고 인수 기준을 붙일 곳이 없다.}
\begin{fieldbox}
\ifwbblank
\g{L1 = [프로젝트] / L2 = [산출물 + 프로젝트 관리] / L3 = [작업패키지, 크기 기준: 스프린트 \hspace{6mm}~\hspace{6mm}회 또는 원가 \hspace{10mm}억]}\par\vspace{4mm}
\g{코드: [1.n.m] / 분해 기준: 산출물 지향(명사), 100\% 규칙, 상호 배타}\par\vspace{4mm}
\g{인프라·공통 기반의 위치: \hrulefill}\par\vspace{4mm}
\g{통합시험·감리·컷오버·이관·종료의 위치: \hrulefill}\par\vspace{4mm}
\g{발주 / 수행 / 공동 담당 구분 기준: \hrulefill}\par\vspace{2mm}
\else
\begin{itemize}
\item L1 = P1. L2 = 헌장 §5 포함 5개 산출물(1.1~1.5) + 프로젝트 관리·통합 검증(1.6). L3 = 작업패키지 43개.
\item 코드 1.n.m. 작업패키지 크기: 스프린트 1~6회 상당(원가 약 0.5~9억). 하한 미만은 상위 패키지에 흡수, 상한 초과는 분할.
\item 분해 기준: 산출물 지향(명사), 100\% 규칙(L2 합 = L1, L3 합 = L2), 상호 배타(같은 인도물이 두 패키지에 없다).
\item 인프라(클라우드·DR·망분리)는 1.5 공통 기반에, 통합시험·성능·보안진단·감리·컷오버 관리는 1.6에 둔다. 컷오버 “매장·주문 1차 / 물류 2차” 분리 옵션(R-01)은 1.3.4와 1.6.9를 분리해 둠으로써 구조에 반영한다.
\item 담당 열의 발주/수행 구분은 FP 12,400점의 안팎을 가른다. 발주 담당 패키지는 수행사 FP에 포함되지 않으며, 발주사 미결정 대기(R-05)는 이 패키지들의 일정 지연으로 나타난다.
\end{itemize}\fi
\end{fieldbox}
\vspace{2mm}
\noindent\textbf{\small 레벨 1·2 트리} \g{(L2 노드에 FP / 원가(억)를 적고 합계가 계획 원가와 같은지 확인)}\par\vspace{1mm}
\begin{center}
\begin{tikzpicture}[
  every node/.style={draw,rounded corners=1pt,align=center,font=\footnotesize,inner sep=2.5pt},
  root/.style={fill=thead,minimum width=70mm,minimum height=9mm},
  l2/.style={minimum width=26mm,minimum height=15mm,text width=25mm}]
\node[root] (r) at (0,0) {\B{\g{1 [프로젝트명]}}{\textbf{1 P1 차세대 주문·재고 통합 플랫폼}}};
\node[l2] (a) at (-7.2,-2.2) {\B{\g{1.1 [산출물 ①]\\[2mm]FP\ \ \ \ \ /\ \ \ \ \ 억}}{1.1 POS·매장 운영\\FP 3,100 / 25.65억}};
\node[l2] (b) at (-4.32,-2.2) {\B{\g{1.2 [산출물 ②]\\[2mm]FP\ \ \ \ \ /\ \ \ \ \ 억}}{1.2 주문 허브\\FP 2,900 / 15.95억}};
\node[l2] (c) at (-1.44,-2.2) {\B{\g{1.3 [산출물 ③]\\[2mm]FP\ \ \ \ \ /\ \ \ \ \ 억}}{1.3 통합 재고 관리\\FP 2,300 / 12.65억}};
\node[l2] (d) at (1.44,-2.2) {\B{\g{1.4 [산출물 ④]\\[2mm]FP\ \ \ \ \ /\ \ \ \ \ 억}}{1.4 가맹 발주 포털\\FP 1,500 / 8.25억}};
\node[l2] (e) at (4.32,-2.2) {\B{\g{1.5 [산출물 ⑤]\\[2mm]FP\ \ \ \ \ /\ \ \ \ \ 억}}{1.5 통합 계층·전환·공통 기반\\FP 2,600 / 64.00억}};
\node[l2] (f) at (7.2,-2.2) {\B{\g{1.6 [프로젝트 관리]\\[2mm]FP — /\ \ \ \ \ 억}}{1.6 프로젝트 관리·통합 검증\\FP — / 19.30억}};
\foreach \x in {a,b,c,d,e,f}{\draw (r.south) -- ++(0,-0.45) -| (\x.north);}
\node[draw=none,font=\footnotesize] at (0,-3.6) {\B{\g{합계 FP [\ \ \ \ \ \ ] / [\ \ \ \ \ \ ]억 = 계획 원가}}{합계 FP 12,400 / 145.80억 = 계획 원가 (⑨ 항목 2)}};
\end{tikzpicture}
\end{center}

% ------------------------------------------------------------------ 8 (가로)
\landscapeon
\secbar{8. WBS 트리 — 레벨 2·3 목록}
\guide{코드, 작업패키지명(명사 — “~규격서”, “~모듈”, “~패키지”), 대표 산출물, 담당(발주/수행/공동, 실명·팀), 릴리스·스프린트 위치, 인수자. 원가는 L2까지만(항목 9), L3 원가는 WBS 사전에. 담당 칸의 발주/수행 구분은 고정가 계약에서 FP 밖의 패키지(조달·단말·정본 판정·협상 지원)를 드러낸다. 작업패키지가 100개를 넘으면 관리 불가, 10개 이하면 원가를 붙일 수 없다. 컷오버·안정화·이관·종료가 없는 WBS는 끝나지 않는다.}
{\scriptsize
\begin{longtable}{|C{7mm}|C{7mm}|C{7mm}|L{12mm}|L{52mm}|L{52mm}|L{44mm}|L{30mm}|L{28mm}|}\hline
\hd{L1} & \hd{L2} & \hd{L3} & \hd{코드} & \hd{작업패키지 (명사)} & \hd{대표 산출물} & \hd{담당 (발주/수행/공동, 실명·팀)} & \hd{릴리스·스프린트} & \hd{인수자} \\ \hline \endhead
\ifwbblank
\rd{7mm}● & & & 1 & \g{[프로젝트]} & & & & \\ \hline
\rd{7mm} & ● & & 1.1 & \g{[L2 산출물 ①]} & & & & \\ \hline
\rd{7mm} & & ● & 1.1.1 & \g{[작업패키지 — 인도물 이름]} & \g{[인도물 목록]} & \g{[발주/수행/공동 — 실명·팀]} & \g{[Rn Sn–Sn]} & \g{[실명]} \\ \hline
\rd{7mm} & & ● & 1.1.2 & & & & & \\ \hline
\rd{7mm} & & ● & 1.1.3 & & & & & \\ \hline
\rd{7mm} & & ● & 1.1.4 & & & & & \\ \hline
\rd{7mm} & & ● & 1.1.5 & & & & & \\ \hline
\rd{7mm} & ● & & 1.2 & \g{[L2 산출물 ②]} & & & & \\ \hline
\rd{7mm} & & ● & 1.2.1 & & & & & \\ \hline
\rd{7mm} & & ● & 1.2.2 & & & & & \\ \hline
\rd{7mm} & & ● & 1.2.3 & & & & & \\ \hline
\rd{7mm} & & ● & 1.2.4 & & & & & \\ \hline
\rd{7mm} & ● & & 1.3 & \g{[L2 산출물 ③]} & & & & \\ \hline
\rd{7mm} & & ● & 1.3.1 & & & & & \\ \hline
\rd{7mm} & & ● & 1.3.2 & & & & & \\ \hline
\rd{7mm} & & ● & 1.3.3 & & & & & \\ \hline
\rd{7mm} & & ● & 1.3.4 & & & & & \\ \hline
\rd{7mm} & ● & & 1.4 & \g{[L2 산출물 ④]} & & & & \\ \hline
\rd{7mm} & & ● & 1.4.1 & & & & & \\ \hline
\rd{7mm} & & ● & 1.4.2 & & & & & \\ \hline
\rd{7mm} & & ● & 1.4.3 & & & & & \\ \hline
\rd{7mm} & ● & & 1.5 & \g{[L2 산출물 ⑤]} & & & & \\ \hline
\rd{7mm} & & ● & 1.5.1 & & & & & \\ \hline
\rd{7mm} & & ● & 1.5.2 & & & & & \\ \hline
\rd{7mm} & & ● & 1.5.3 & & & & & \\ \hline
\rd{7mm} & & ● & 1.5.4 & & & & & \\ \hline
\rd{7mm} & & ● & 1.5.5 & & & & & \\ \hline
\rd{7mm} & ● & & 1.6 & \g{[프로젝트 관리·통합 검증]} & & & & \\ \hline
\rd{7mm} & & ● & 1.6.1 & \g{[기획·기준선]} & & & & \\ \hline
\rd{7mm} & & ● & 1.6.2 & \g{[PMO·형상·품질]} & & & & \\ \hline
\rd{7mm} & & ● & 1.6.3 & \g{[통합시험·UAT]} & & & & \\ \hline
\rd{7mm} & & ● & 1.6.4 & \g{[컷오버·안정화]} & & & & \\ \hline
\rd{7mm} & & ● & 1.6.5 & \g{[이관·종료]} & & & & \\ \hline
\rd{7mm} & & ● & & \g{(필요 시 행 추가 — 별지)} & & & & \\ \hline
\else
● & & & 1 & \textbf{P1 차세대 주문·재고 통합 플랫폼} & & & & \\ \hline
& ● & & \textbf{1.1} & \textbf{POS·매장 운영 시스템} & & & & \\ \hline
& & ● & 1.1.1 & POS 요구·화면 설계서 & 화면 정의서, 거래 시나리오 & 수행 POS팀 + 오정근(요구) & R0 S1–S6 & 오정근 \\ \hline
& & ● & 1.1.2 & 매장 거래 처리 모듈 & 판매·결제·반품·할인·멤버십 & 수행 POS팀 & R1–R2 S7–S18 & 박세연 \\ \hline
& & ● & 1.1.3 & 매장 운영 모듈 & 개점·마감·실사·폐기·발주 연계 & 수행 POS팀 & R2–R3 S13–S24 & 오정근 \\ \hline
& & ● & 1.1.4 & 결제·PG·밴 인터페이스 7본 & 인터페이스 7본, 문서 2본 복원 & 수행 인터페이스팀 & R2 S13–S18 & 박세연 \\ \hline
& & ● & 1.1.5 & 직영 단말 477대 조달·설치 & 단말·주변기기, 설치 확인서 & 발주 (강현도 조달, 매장운영 설치) & 2026-12~2027-02 & 오정근 \\ \hline
& & ● & 1.1.6 & 가맹 단말 호환성 조사·SW 배포 패키지 & 388점 기종 조사서, 비호환 목록, 배포 패키지 & 발주 (오정근 조사) + 수행 POS팀 & 조사 2026-04, 패키지 R4 & 오정근 \\ \hline
& & ● & 1.1.7 & 파일럿 12개점 배포·임시 연동 & 파일럿 배포 기록, 임시 재고 동기화 & 공동 (P3 PM, 수행 POS팀) & R4 S27–S28 & 오정근, P3 PM \\ \hline
& & ● & 1.1.8 & 컷오버 펌웨어 일괄배포·롤백 패키지 & 배포 패키지, 롤백 절차서, 리허설 기록 & 수행 POS팀 + 발주 운영 & R4 S29–S30 & 박세연, 박준영 \\ \hline
& ● & & \textbf{1.2} & \textbf{주문 허브} & & & & \\ \hline
& & ● & 1.2.1 & 주문 원장 모델·채널 설계서 & 원장 모델, 채널별 매핑 규격 & 수행 주문팀 + 오정근 & R0 S1–S6 & 박세연 \\ \hline
& & ● & 1.2.2 & 주문 원장 코어 모듈 & 주문 원장, 상태 관리, 정합 검증 & 수행 주문팀 & R3 S19–S24 & 오정근 \\ \hline
& & ● & 1.2.3 & 자사앱·배달3사 채널 연동 & 배달3사 3본 재연결, 앱 API 규격·연동 & 수행 인터페이스팀 (앱측: 벤더) & R3 S19–S24 & 오정근, 박세연 \\ \hline
& & ● & 1.2.4 & B2B EDI 채널 연동 & 편의점·식품사·원료유통 EDI 접수 & 수행 인터페이스팀 & R4 S25–S26 & 조성태(공장), 오정근 \\ \hline
& & ● & 1.2.5 & 옴니채널 시나리오 (기본 2종 / 고도화 이연) & 픽업·매장 간 이동 기본 / 자동 제안·예약·BI 2종 & 수행 주문팀 & 기본 R4 S25–S26 / 고도화 R5 S32–S34 & 오정근 \\ \hline
& ● & & \textbf{1.3} & \textbf{통합 재고 관리} & & & & \\ \hline
& & ● & 1.3.1 & 재고 원장 모델·규격서 (P2·3PL 협상 지원 자료 포함) & 원장 모델, P2 인터페이스 규격(T0 전), 3PL 부하 산정서 & 발주 (박세연, 한동석 협의) + 수행 재고팀 & R0 S1–S6, T0 2026-04-24 & 박세연, 프로그램 PMO장 \\ \hline
& & ● & 1.3.2 & 재고 원장 코어 모듈 & 3개 재고원 단일 원장, 15분 반영 & 수행 재고팀 & R2 S13–S18 & 한동석 \\ \hline
& & ● & 1.3.3 & 3PL WMS 준실시간 API·폴백 & 연동 명세서, P1측 API, 야간 배치 2회 폴백, 자동 전환 & 수행 인터페이스팀 (협상: 한동석) & R2 S13–S16 & 한동석, 박세연 \\ \hline
& & ● & 1.3.4 & P2 설비 인터페이스 P1측 수신 구현·시험 & 수신 모듈, 통합시험 결과(MC 후) & 수행 인터페이스팀 (설비측: P2) & R3 구현, 시험 2027-02-12~18 & 프로그램 PMO장 \\ \hline
& & ● & 1.3.5 & 공장 출하 재고 연동 (별도 법인) & 온담푸드시스템 연동 2본, 계약·보안 정책 반영 & 수행 인터페이스팀 (계약: 최영주·조성태) & R3 S19–S24 & 조성태, 박세연 \\ \hline
& ● & & \textbf{1.4} & \textbf{가맹 발주 포털} & & & & \\ \hline
& & ● & 1.4.1 & 포털 요구·UX 설계서 (협의회 참여) & 화면 정의서, 점주 인터뷰 기록 & 수행 포털팀 + 오정근·서정필 & R0 S1–S6 & 오정근 \\ \hline
& & ● & 1.4.2 & 발주·마감·출고 확인·배송 추적 모듈 & 포털 모듈, 엑셀 업로드·모바일 & 수행 포털팀 & R1–R3 S7–S22 & 오정근 \\ \hline
& & ● & 1.4.3 & 프로토타입 시연·선행 릴리스 & 프로토타입(2026-05), 선행 릴리스(2026-Q4) & 수행 포털팀 + 발주 매장운영 & R1 S9, R3 S23–S24 & 오정근, 서정필 \\ \hline
& & ● & 1.4.4 & 388점 계정·온보딩 & 계정 발급, 온보딩 안내(교육은 P3) & 발주 매장운영 + 수행 포털팀 & R3–R4 & 오정근 \\ \hline
& ● & & \textbf{1.5} & \textbf{통합 계층·데이터 전환·공통 기반} & & & & \\ \hline
& & ● & 1.5.1 & 아키텍처·비기능 설계서 & 아키텍처 정의서, 가용성 99.5\%·망분리·규제 요건 매핑 & 수행 아키텍트 + 박세연·최영주 & R0 S1–S6 & 박세연 \\ \hline
& & ● & 1.5.2 & API 게이트웨이 구축 & 상용 API GW 설치·정책, 47본 라우팅 & 수행 통합팀 + SW 벤더 & R1 S7–S12 & 박세연 \\ \hline
& & ● & 1.5.3 & 상품·고객 MDM 구축 & MDM 설치, 마스터 모델, 정제 규칙 & 수행 통합팀 + SW 벤더 & R1–R2 S7–S18 & 박세연 \\ \hline
& & ● & 1.5.4 & SAP 애드온·SAP 인터페이스 15본 & 애드온, OD-POS↔SAP 11본·가맹발주↔SAP 4본 대체 & 수행 SAP팀 & R2–R3 S13–S22 & 박세연 \\ \hline
& & ● & 1.5.5 & 인터페이스 47본 재구축·설계문서 & 47본 설계문서(13본 복원), 재연결 & 수행 인터페이스팀 + 박세연(복원 지식) & 복원 S2–S6, 재구축 R1–R4 & 박세연 \\ \hline
& & ● & 1.5.6 & 클라우드·DR·망분리·APM 인프라 & 클라우드 환경, DR, 망분리·보안장비, APM & 발주 (박세연, 강현도 조달) + 벤더 & 조달 2026-02~04, DR 2026-12~2027-01 & 박세연 \\ \hline
& & ● & 1.5.7 & 데이터 전환 설계·매핑 (1,842 테이블) & 매핑 정의서, 전환 프로그램, 아카이브 정책 & 데이터 전환 용역사 + 박세연 & 2026-03-30~10-23 & 박세연 \\ \hline
& & ● & 1.5.8 & 마스터 6종 정제·정본 판정 (G2) & 정제 결과, 정본 판정서, 동결 규칙 & 발주 (박세연 R, 감리 확인) + 용역사 & 2026-10-26~11-20 → G2 & 프로그램 이사회 \\ \hline
& & ● & 1.5.9 & 전환 리허설 1~3회차 & 리허설 결과 3건, 검증 6종 보고서 & 용역사 + 수행 + 발주 운영 & 2026-11-23~2027-01-22 & 박세연 \\ \hline
& & ● & 1.5.10 & 본 전환 실행 & 전환 완료 보고, 검증 6종 & 용역사 + 수행 & 컷오버 2027-02-26~28 & 박세연 \\ \hline
& ● & & \textbf{1.6} & \textbf{프로젝트 관리·통합 검증} & & & & \\ \hline
& & ● & 1.6.1 & 기획·기준선 (G1) & 범위·일정·원가 기준선, FP 확인서 BL-01 & 발주 P1 PM + 이재현 & 2026-03-30~04-24 → G1 & 프로그램 이사회 \\ \hline
& & ● & 1.6.2 & PMO 운영·형상관리·품질보증 & 월 보고, 형상 저장소, 품질 지표 & 발주 P1 PMO 3명 & 상시 & P1 PM \\ \hline
& & ● & 1.6.3 & 요구사항 관리 (RTM·월 동결) & RTM, 동결 기록 12회, 기준선 재확인서 & 발주 PMO + 수행 & 상시 & P1 PM, 오정근 \\ \hline
& & ● & 1.6.4 & 통합시험 & TC 400건, 결과 보고 & 수행 시험팀 + 발주 검증 & S25–S27 (2026-12-07~2027-01-15) & 박세연 \\ \hline
& & ● & 1.6.5 & 성능·부하시험 & 부하시험 보고(피크 5,100 케이스/h 기준 재고 API) & 3자 시험 + 수행 & 2027-01 & 박세연, 박준영 \\ \hline
& & ● & 1.6.6 & 보안취약점 진단 & 진단 보고 2회(규제 전·컷오버 전), High 0건 & 3자 진단 & 2026-08, 2027-01 & 최영주, 박세연 \\ \hline
& & ● & 1.6.7 & UAT & UAT TC 360건, 종료 기준 판정 & 발주 현업 + 수행 & S28–S30 (2027-01-18~02-18) & 오정근 \\ \hline
& & ● & 1.6.8 & 외부 감리 대응 (G1·G2·G3·G4) & 감리 지적·시정 기록 & 발주 PMO (정하늘 감리) & 게이트별 & P1 PM \\ \hline
& & ● & 1.6.9 & 컷오버 계획·실행·안정화 (매장·주문 / 물류 분리 옵션) & 컷오버 계획, go/no-go 자료, 하이퍼케어 보고 & 공동 (P1 PM, 이재현, 박준영) & G3 → 2027-02-26~28 → S31–S34 & 정미란(go/no-go) \\ \hline
& & ● & 1.6.10 & 운영 이관 기술문서·소스 인수 (G4) & 기술 문서, 클린 빌드 재현, SBOM & 발주 (박세연, 형상 담당) + 수행 & S35–S36 & 박준영, 박세연 \\ \hline
& & ● & 1.6.11 & 종료 (보고서·교훈) & 종료 보고서, 교훈 기록 & 발주 P1 PM & G4 & 정미란 \\ \hline
\fi
\end{longtable}}
\B{}{\small 작업패키지 43개 = 발주 담당 10(1.1.5, 1.1.6, 1.3.1, 1.5.6, 1.5.8, 1.6.1~1.6.3, 1.6.8, 1.6.10) + 수행 담당 28 + 공동 5.\par}

% ------------------------------------------------------------------ 9 (가로)
\secbar[60mm]{9. L2 원가·FP 배분}
\guide{행 = L2, 열 = 원가 항목, 합계 행·열. SI는 FP 배분 × 단가, 나머지는 항목의 성격으로 배분한다. \textbf{합계가 원가 기준선(⑨)의 계획 원가와 정확히 같아야 한다} — 다르면 WBS 또는 원가 기준선 중 하나가 100\% 규칙을 어긴 것이다. SI만 배분하고 라이선스·인프라·PMO를 “공통”으로 남기면 L2 합이 기준선과 맞지 않는다.}
{\footnotesize
\begin{tabularx}{\linewidth}{|Y|C{18mm}|C{24mm}|C{18mm}|C{18mm}|C{18mm}|C{20mm}|C{20mm}|C{18mm}|C{22mm}|}\hline
\hd{L2} & \hd{FP} & \hd{SI (FP × 단가)} & \hd{상용SW} & \hd{인프라} & \hd{단말} & \hd{데이터 전환} & \hd{시험·감리} & \hd{PMO} & \hd{합계 (억)} \\ \hline
\ifwbblank
\rd{8mm}1.1 \g{[\ ]} & & & & & & & & & \\ \hline
\rd{8mm}1.2 & & & & & & & & & \\ \hline
\rd{8mm}1.3 & & & & & & & & & \\ \hline
\rd{8mm}1.4 & & & & & & & & & \\ \hline
\rd{8mm}1.5 & & & & & & & & & \\ \hline
\rd{8mm}1.6 \g{[프로젝트 관리]} & \g{—} & \g{—} & & & & & & & \\ \hline
\rd{8mm}\textbf{합계} & \g{[= 총 FP]} & & & & & & & & \g{[= 계획 원가]} \\ \hline
\else
1.1 POS·매장 운영 & 3,100 & 17.05 & — & — & 8.60 & — & — & — & \textbf{25.65} \\ \hline
1.2 주문 허브 & 2,900 & 15.95 & — & — & — & — & — & — & \textbf{15.95} \\ \hline
1.3 통합 재고 & 2,300 & 12.65 & — & — & — & — & — & — & \textbf{12.65} \\ \hline
1.4 가맹 발주 포털 & 1,500 & 8.25 & — & — & — & — & — & — & \textbf{8.25} \\ \hline
1.5 통합 계층·전환·공통 기반 & 2,600 & 14.30 & 19.40 & 20.90 & — & 9.40 & — & — & \textbf{64.00} \\ \hline
1.6 프로젝트 관리·통합 검증 & — & — & — & — & — & — & 9.70 & 9.60 & \textbf{19.30} \\ \hline
\textbf{합계} & \textbf{12,400} & \textbf{68.20} & \textbf{19.40} & \textbf{20.90} & \textbf{8.60} & \textbf{9.40} & \textbf{9.70} & \textbf{9.60} & \textbf{145.80} \\ \hline
\fi
\end{tabularx}}
\B{}{\small SI 열 = FP × 55만 원. 상용SW 19.40·인프라 20.90·전환 9.40·PMO 9.60은 승인 내역(23.4 / 23.9 / 11.4 / 12.6)에서 조정된 기준선 값 — 조정 근거는 ⑨ 항목 3. 합계 145.80 = 계획 원가.\par}
\landscapeoff

\begin{center}\bfseries\large WBS 사전 (작업패키지 카드)\end{center}
% ------------------------------------------------------------------ 10
\secbar[60mm]{10. WBS 사전 — 작업패키지 카드}
\guide{작업패키지마다 한 장. 코드·명칭 / 산출물(인도물 목록) / 인수 기준과 판정 방법 / 담당·인수자 / 선행·후행·외부 의존(상대 실명·기한) / 추정(값과 근거를 같이: “FP 420 × 55만 원 = 2.31억, 근거: FP 산정서 v0.9 §3.3”) / 관련 요구(RTM ID)·리스크(R-ID) / 가정. PRINCE2의 Product Description과 같은 항목. 43장을 다 쓰려다 아무것도 쓰지 않는다 — \textbf{외부 의존이 있는 것, 임계경로에 있는 것, 발주사 담당인 것부터.}}
% \wbscard{코드/명칭}{산출물}{인수 기준/판정}{담당/인수자}{선행/후행/외부}{추정}{요구/리스크}{가정}
\newcommand{\wbscard}[8]{%
{\footnotesize
\begin{tabularx}{\linewidth}{|L{34mm}|Y|}\hline
\hd{항목} & \hd{내용} \\ \hline
코드 / 명칭 & #1 \\ \hline
산출물 (인도물) & #2 \\ \hline
인수 기준 / 판정 방법 & #3 \\ \hline
담당 / 인수자 & #4 \\ \hline
선행 / 후행 / 외부 의존 & #5 \\ \hline
추정 (값 + 근거) & #6 \\ \hline
관련 요구 / 리스크 & #7 \\ \hline
가정 & #8 \\ \hline
\end{tabularx}}\par\vspace{3mm}}
\ifwbblank
\foreach \i in {1,2,3}{%
\needspace{95mm}\noindent\textbf{카드 \i} \g{— 코드 \hspace{20mm}}\par\vspace{1mm}
\wbscard{\rd{9mm}\g{[1.n.m] / [명칭 — 명사]}}
 {\rd{16mm}\g{① [\ ] ② [\ ] ③ [\ ]}}
 {\rd{16mm}\g{(1) [판정 가능 문장] (2) [\ ] / [시험 ID·검수·감리·문서 확인]}}
 {\rd{9mm}\g{[실명·팀] / [실명]}}
 {\rd{12mm}\g{[코드] / [코드] / [상대 실명, 기한, 미이행 시 조치]}}
 {\rd{12mm}\g{FP [\ ] × 단가 = [\ ]억 / 공수 [\ ] 인월 / 기간 [\ ]영업일 (Sn–Sn) / 근거 [\ ]}}
 {\rd{9mm}\g{[StRS-nnn, HR-nn] / [R-nn (소유자)]}}
 {\rd{12mm}\g{[가정] (확인: 실명, 시점) — 틀리면 [영향]}}}
\else
\needspace{80mm}\noindent\textbf{카드 1.3.3 — 3PL WMS 준실시간 API·폴백}\par\vspace{1mm}
\wbscard{1.3.3 / 3PL 대성로지스 WMS 준실시간 연동 API와 야간 배치 폴백}
 {① 연동 명세서 v1.0(항목·주기·오류 처리·부하) ② 부하 산정서(한동석 협상용, 2026-03) ③ P1측 수신·송신 API ④ 야간 배치 2회(21:00·03:00) 폴백 모듈 ⑤ 준실시간↔폴백 자동 전환 기능과 전환 시험 결과}
 {① 준실시간 모드: 센터 입출고가 재고 원장에 15분 이내 반영(TC-INV-021~030) ② 폴백 모드: 배치 2회 자동 실행·실패 시 재시도·알림(TC-INV-031~036) ③ 모드 전환 시 재고 원장 정합 100\%(TC-INV-037~040) ④ 부하 산정서에 대성 서면 회신 첨부 / 통합시험, 성능시험}
 {수행 인터페이스팀(리더: 한빛디지털 인터페이스 PL) + 발주 물류본부 센터운영 담당 / 한동석 물류본부장, 박세연 IT본부장}
 {선행 1.3.1(규격), 1.3.2(원장 코어) / 후행 1.6.4(통합시험), 1.6.5(성능) / \textbf{외부: 대성로지스 계약 변경 타결 — 한동석, 2026-06-30. 미타결 시 ④·⑤만 인도하고 ③은 명세·구현 완료 상태로 보관}}
 {FP 420 × 55만 원 = 2.31억 (근거: FP 산정서 v0.9 §3.3 — 인터페이스 5본 × 트랜잭션·데이터 기능) / 공수 4.2 인월 / 기간 40영업일 (S13–S16, 2026-06-22~08-14) / 발주사 공수 0.5 인월(부하 산정·검증)}
 {StRS-031, 032, 034 (HR-01) / R-03 (소유자 한동석), R-12(정확도 분모)}
 {대성 WMS의 기존 EDI 형식이 준실시간 API에도 유지된다 — 아니면 매핑 재작업 FP +60 (확인: 수행 인터페이스 PL, 2026-04)}
\needspace{80mm}\noindent\textbf{카드 1.5.8 — 마스터 6종 정제·정본 판정 (G2)}\par\vspace{1mm}
\wbscard{1.5.8 / 마스터 데이터 6종(상품·고객·매장·거래처·가격·재고) 정제 및 정본 판정}
 {① 정제 규칙서(중복·결측·코드 체계) ② 정제 결과 보고(6종 건수·중복률 전후) ③ 정본 판정서(G2 2026-11-27 프로그램 이사회) ④ 정본 이후 변경 동결 규칙(CCB 의결, R-10 회피) ⑤ 610개 종속 테이블 매핑 확정 목록}
 {① 6종 중복률 0.5\% 이하, 필수 항목 결측 0 ② 상품 코드 체계 확정서에 오정근·조성태·강현도 서명 ③ 감리(정하늘) 확인 의견 첨부 — 박세연은 R이며 A가 아님(RACI 23행, 이해 충돌 기록) ④ 동결 규칙 CCB 의결서 / 감리 문서 검토, 정제 결과 표본 대사}
 {발주 IT본부(박세연 R) + 데이터 전환 용역사 + 현업 데이터 오너(상품: 오정근·조성태 / 거래처·가격: 강현도 / 매장: 매장운영 / 고객: 영업본부 CRM) / 프로그램 이사회(A)}
 {선행 1.5.7(매핑), 1.5.3(MDM) / 후행 1.5.9(리허설), 1.2.2·1.3.2(원장 코어의 마스터 참조) / 외부: 신규 브랜드·B2B 상품 코드 요구 — 영업본부, 2026-10까지 접수 마감}
 {데이터 전환 용역 9.4억 중 2.6억 [추가 설정] / 발주 공수 3.0 인월(데이터 오너 검토) / 기간 20영업일 (2026-10-26~11-20), G2 여유 5영업일 / \textbf{비가역 — 판정 후 되돌림 3.8억·6주(R-10)}}
 {StRS-141, 161 (HR-07, HR-08), SyRS 데이터 요구 42건 / R-10 (소유자 박세연), R-06 (원 개발자 지식)}
 {온담푸드시스템 별도 법인의 상품 마스터가 그룹 코드 체계를 수용한다 (확인: 조성태, 2026-08)}
\needspace{80mm}\noindent\textbf{카드 1.1.6 — 가맹 단말 호환성 조사·SW 배포 패키지}\par\vspace{1mm}
\wbscard{1.1.6 / 가맹 388점 POS 단말 호환성 전수 조사와 가맹용 SW 배포 패키지}
 {① 388점 단말 기종·OS·메모리 조사서(회수율 100\%, 2026-04-30) ② 비호환 기종 목록과 매장 수, 대응안 3개(교체·임대·수기 병행) ③ 가맹용 SW 배포 패키지(기종별 2~3종) ④ 파일럿 12개점 중 구형 기종 매장 2개 이상 포함 배포 결과}
 {① 조사 회수율 100\%, 기종 분류 완료 ② 비호환 매장 비율 확정(10\% 초과 시 R-13 트리거) ③ 배포 패키지가 호환 기종 전 유형에서 설치·거래 처리 성공(TC-POS-150~180) ④ 파일럿 구형 기종 배포 성공률 100\% / 조사 보고, UAT, 파일럿 기록}
 {발주 영업본부 매장운영(조사, 오정근 A) + 수행 POS팀(패키지) / 오정근}
 {선행 1.1.1 / 후행 1.1.7(파일럿), 1.1.8(일괄배포) / \textbf{외부: 가맹 부속서 개정 388점 동의 — 오정근, 2026-08-14. 교체 비용 부담 결정은 P1 범위 밖(헌장 §5 제외)}}
 {FP 180 × 55만 원 = 0.99억(패키지) + 조사 비용 0.3억 [추가 설정, 발주 매장운영 출장·설문] / 발주 공수 2.0 인월 / 기간: 조사 40영업일(2026-03~04), 패키지 S25–S28}
 {StRS-071, 091 (HR-03, HR-04) / R-02 (소유자 오정근), R-13 (소유자 P1 PM)}
 {2013년 이전 기종은 약 90대로 추정되며(프리모템 원문 11) 조사로 확정한다. 비호환 매장의 수기 병행은 최대 4주 (확인: 오정근, 2026-05)}
\small 나머지 40개 카드는 G1(2026-04-30)까지 같은 양식으로 작성하고 스프린트마다 갱신한다 (작성 순서: 외부 의존 → 임계경로 → 발주 담당).
\fi

\end{document}
